%=============================================================================
\section{Appendix: Plug-and-Play C++ Modules}
%=============================================================================

This appendix provides production-ready, header-only C++ implementations of every major algorithm discussed in this note. Each module is self-contained --- just drop the header into your project and use it. All modules are designed for embedded systems (no dynamic allocation, no exceptions, no STL containers).

\textbf{Design philosophy:}
\begin{itemize}
    \item \textbf{Header-only}: No \texttt{.cpp} files needed. Include and go.
    \item \textbf{No heap allocation}: Everything lives on the stack or is statically sized via templates.
    \item \textbf{Minimal dependencies}: Only \texttt{<cmath>} and \texttt{<cstring>}. No STL.
    \item \textbf{ISR-safe}: All \texttt{update()} methods are safe to call from timer interrupts.
    \item \textbf{Template-based}: Matrix dimensions are compile-time constants for zero-overhead abstraction.
\end{itemize}

\subsection{Low-Pass Filter}

First-order IIR low-pass filter. One line to construct, one line to use.

\begin{verbatim}
// lpf.hpp
#pragma once
#include <cmath>

class LowPassFilter {
public:
    LowPassFilter() : alpha_(1.0f), y_prev_(0.0f), initialized_(false) {}

    /// Configure with cutoff frequency (Hz) and sampling period (s)
    void configure(float cutoff_hz, float dt) {
        float rc = 1.0f / (2.0f * M_PI * cutoff_hz);
        alpha_ = dt / (rc + dt);
    }

    /// Configure directly with smoothing factor (0 < alpha <= 1)
    void set_alpha(float alpha) { alpha_ = alpha; }

    /// Process one sample. Call this at your fixed sample rate.
    float update(float x) {
        if (!initialized_) {
            y_prev_ = x;
            initialized_ = true;
            return x;
        }
        y_prev_ = alpha_ * x + (1.0f - alpha_) * y_prev_;
        return y_prev_;
    }

    /// Reset filter state
    void reset() { initialized_ = false; y_prev_ = 0.0f; }

    /// Get current filtered value without feeding new input
    float value() const { return y_prev_; }

private:
    float alpha_;
    float y_prev_;
    bool  initialized_;
};
\end{verbatim}

\textbf{Usage:}
\begin{verbatim}
    LowPassFilter gyro_filter;
    gyro_filter.configure(30.0f, 0.001f);  // 30 Hz cutoff, 1 kHz loop

    // In your control loop ISR:
    float gyro_clean = gyro_filter.update(gyro_raw);
\end{verbatim}

\subsection{PID Controller}

Full-featured PID with anti-windup (back-calculation), derivative-on-measurement, filtered derivative, and output clamping.

\begin{verbatim}
// pid.hpp
#pragma once
#include <cmath>

struct PIDConfig {
    float kp = 0.0f;
    float ki = 0.0f;
    float kd = 0.0f;
    float dt = 0.001f;         // Sampling period (s)
    float out_min = -1e6f;     // Output lower limit
    float out_max =  1e6f;     // Output upper limit
    float integral_max = 1e6f; // Integral term clamp
    float d_filter_N = 10.0f;  // Derivative filter coefficient (higher = less filtering)
};

class PID {
public:
    PID() = default;
    explicit PID(const PIDConfig& cfg) : cfg_(cfg) {}

    void configure(const PIDConfig& cfg) { cfg_ = cfg; }

    /// Compute PID output. Call once per control cycle.
    /// setpoint: desired value, measurement: actual value
    float update(float setpoint, float measurement) {
        float error = setpoint - measurement;

        // --- Proportional ---
        float p_term = cfg_.kp * error;

        // --- Integral with back-calculation anti-windup ---
        integral_ += cfg_.ki * error * cfg_.dt + anti_windup_;
        integral_ = clamp(integral_, -cfg_.integral_max, cfg_.integral_max);

        // --- Derivative on measurement (no derivative kick) ---
        //     with first-order low-pass filter
        float d_raw = -(measurement - meas_prev_) / cfg_.dt;
        float alpha_d = cfg_.dt * cfg_.d_filter_N
                      / (1.0f + cfg_.dt * cfg_.d_filter_N);
        d_filtered_ = alpha_d * d_raw + (1.0f - alpha_d) * d_filtered_;
        float d_term = cfg_.kd * d_filtered_;

        meas_prev_ = measurement;

        // --- Total output ---
        float output_raw = p_term + integral_ + d_term;
        float output = clamp(output_raw, cfg_.out_min, cfg_.out_max);

        // --- Back-calculation anti-windup ---
        if (cfg_.ki != 0.0f) {
            anti_windup_ = (output - output_raw) * (1.0f / cfg_.ki)
                         * cfg_.dt * 0.5f;
        } else {
            anti_windup_ = 0.0f;
        }

        return output;
    }

    void reset() {
        integral_ = 0.0f;
        d_filtered_ = 0.0f;
        meas_prev_ = 0.0f;
        anti_windup_ = 0.0f;
    }

    // Access internal state for debugging
    float get_integral() const { return integral_; }

private:
    PIDConfig cfg_;
    float integral_    = 0.0f;
    float d_filtered_  = 0.0f;
    float meas_prev_   = 0.0f;
    float anti_windup_ = 0.0f;

    static float clamp(float v, float lo, float hi) {
        return v < lo ? lo : (v > hi ? hi : v);
    }
};
\end{verbatim}

\textbf{Usage:}
\begin{verbatim}
    PIDConfig cfg;
    cfg.kp = 6.0f;  cfg.ki = 0.3f;  cfg.kd = 3.0f;
    cfg.dt = 0.001f;
    cfg.out_min = -30000.0f;  cfg.out_max = 30000.0f;  // GM6020 range
    cfg.integral_max = 10000.0f;

    PID yaw_pid(cfg);

    // In ISR:
    float torque_cmd = yaw_pid.update(target_angle, current_angle);
\end{verbatim}

\subsection{Setpoint Ramp (LPF Soft-Start)}

A dedicated setpoint smoother --- wraps LPF to give ramp-limited setpoint transitions.

\begin{verbatim}
// ramp.hpp
#pragma once
#include "lpf.hpp"

class SetpointRamp {
public:
    /// ramp_time: approximate time to reach 95% of a step change (in seconds)
    void configure(float ramp_time, float dt) {
        float cutoff = 3.0f / (2.0f * M_PI * ramp_time);  // 3*tau ~ 95%
        filter_.configure(cutoff, dt);
    }

    float update(float raw_setpoint) {
        return filter_.update(raw_setpoint);
    }

    void reset() { filter_.reset(); }

private:
    LowPassFilter filter_;
};
\end{verbatim}

\subsection{Complementary Filter (IMU Tilt)}

Fuses accelerometer (tilt from gravity) and gyroscope (angular rate).

\begin{verbatim}
// complementary_filter.hpp
#pragma once
#include <cmath>

class ComplementaryFilter {
public:
    /// alpha: gyro trust factor (0.95-0.99 typical)
    void configure(float alpha, float dt) {
        alpha_ = alpha;
        dt_ = dt;
    }

    /// accel_angle: tilt from accelerometer (rad)
    /// gyro_rate:   angular rate from gyro (rad/s)
    /// returns:     fused angle estimate (rad)
    float update(float accel_angle, float gyro_rate) {
        angle_ = alpha_ * (angle_ + gyro_rate * dt_)
               + (1.0f - alpha_) * accel_angle;
        return angle_;
    }

    void reset(float initial_angle = 0.0f) { angle_ = initial_angle; }
    float angle() const { return angle_; }

private:
    float alpha_ = 0.98f;
    float dt_    = 0.001f;
    float angle_ = 0.0f;
};
\end{verbatim}

\subsection{Lightweight Matrix Library}

Compile-time sized matrix for embedded Kalman/LQR. No heap, no STL.

\begin{verbatim}
// mat.hpp
#pragma once
#include <cstring>
#include <cmath>

template<int ROWS, int COLS>
struct Mat {
    float data[ROWS][COLS] = {};

    float& operator()(int r, int c)       { return data[r][c]; }
    float  operator()(int r, int c) const { return data[r][c]; }

    static Mat zeros() { Mat m; return m; }

    static Mat identity() {
        static_assert(ROWS == COLS, "Identity requires square matrix");
        Mat m;
        for (int i = 0; i < ROWS; ++i) m.data[i][i] = 1.0f;
        return m;
    }

    // Matrix addition
    Mat operator+(const Mat& rhs) const {
        Mat r;
        for (int i = 0; i < ROWS; ++i)
            for (int j = 0; j < COLS; ++j)
                r.data[i][j] = data[i][j] + rhs.data[i][j];
        return r;
    }

    // Matrix subtraction
    Mat operator-(const Mat& rhs) const {
        Mat r;
        for (int i = 0; i < ROWS; ++i)
            for (int j = 0; j < COLS; ++j)
                r.data[i][j] = data[i][j] - rhs.data[i][j];
        return r;
    }

    // Matrix multiplication
    template<int COLS2>
    Mat<ROWS, COLS2> operator*(const Mat<COLS, COLS2>& rhs) const {
        Mat<ROWS, COLS2> r;
        for (int i = 0; i < ROWS; ++i)
            for (int j = 0; j < COLS2; ++j)
                for (int k = 0; k < COLS; ++k)
                    r.data[i][j] += data[i][k] * rhs.data[k][j];
        return r;
    }

    // Scalar multiplication
    Mat operator*(float s) const {
        Mat r;
        for (int i = 0; i < ROWS; ++i)
            for (int j = 0; j < COLS; ++j)
                r.data[i][j] = data[i][j] * s;
        return r;
    }

    // Transpose
    Mat<COLS, ROWS> T() const {
        Mat<COLS, ROWS> r;
        for (int i = 0; i < ROWS; ++i)
            for (int j = 0; j < COLS; ++j)
                r.data[j][i] = data[i][j];
        return r;
    }

    // 2x2 or 3x3 inverse (small matrix only, used by Kalman)
    Mat inverse() const;
};

// 1x1 inverse
template<> inline Mat<1,1> Mat<1,1>::inverse() const {
    Mat<1,1> r;
    r.data[0][0] = 1.0f / data[0][0];
    return r;
}

// 2x2 inverse
template<> inline Mat<2,2> Mat<2,2>::inverse() const {
    Mat<2,2> r;
    float det = data[0][0]*data[1][1] - data[0][1]*data[1][0];
    float inv_det = 1.0f / det;
    r.data[0][0] =  data[1][1] * inv_det;
    r.data[0][1] = -data[0][1] * inv_det;
    r.data[1][0] = -data[1][0] * inv_det;
    r.data[1][1] =  data[0][0] * inv_det;
    return r;
}

// 3x3 inverse
template<> inline Mat<3,3> Mat<3,3>::inverse() const {
    Mat<3,3> r;
    float det = data[0][0]*(data[1][1]*data[2][2]-data[1][2]*data[2][1])
              - data[0][1]*(data[1][0]*data[2][2]-data[1][2]*data[2][0])
              + data[0][2]*(data[1][0]*data[2][1]-data[1][1]*data[2][0]);
    float inv = 1.0f / det;
    r(0,0)=(data[1][1]*data[2][2]-data[1][2]*data[2][1])*inv;
    r(0,1)=(data[0][2]*data[2][1]-data[0][1]*data[2][2])*inv;
    r(0,2)=(data[0][1]*data[1][2]-data[0][2]*data[1][1])*inv;
    r(1,0)=(data[1][2]*data[2][0]-data[1][0]*data[2][2])*inv;
    r(1,1)=(data[0][0]*data[2][2]-data[0][2]*data[2][0])*inv;
    r(1,2)=(data[0][2]*data[1][0]-data[0][0]*data[1][2])*inv;
    r(2,0)=(data[1][0]*data[2][1]-data[1][1]*data[2][0])*inv;
    r(2,1)=(data[0][1]*data[2][0]-data[0][0]*data[2][1])*inv;
    r(2,2)=(data[0][0]*data[1][1]-data[0][1]*data[1][0])*inv;
    return r;
}
\end{verbatim}

\subsection{Kalman Filter}

Generic discrete-time Kalman filter. Template parameters set the state and measurement dimensions at compile time.

\begin{verbatim}
// kalman.hpp
#pragma once
#include "mat.hpp"

template<int N_STATE, int N_MEAS, int N_INPUT = 1>
class KalmanFilter {
public:
    using StateVec  = Mat<N_STATE, 1>;
    using MeasVec   = Mat<N_MEAS, 1>;
    using InputVec  = Mat<N_INPUT, 1>;
    using StateMat  = Mat<N_STATE, N_STATE>;
    using InputMat  = Mat<N_STATE, N_INPUT>;
    using MeasMat   = Mat<N_MEAS, N_STATE>;
    using GainMat   = Mat<N_STATE, N_MEAS>;
    using MeasCov   = Mat<N_MEAS, N_MEAS>;

    StateMat A;   // State transition
    InputMat B;   // Input matrix
    MeasMat  C;   // Measurement matrix (called H in some texts)
    StateMat Q;   // Process noise covariance
    MeasCov  R;   // Measurement noise covariance

    StateVec x;   // State estimate
    StateMat P;   // Estimate covariance

    /// Predict step: propagate state and covariance forward
    void predict(const InputVec& u) {
        x = A * x + B * u;
        P = A * P * A.T() + Q;
    }

    /// Predict step (no input)
    void predict() {
        x = A * x;
        P = A * P * A.T() + Q;
    }

    /// Update step: incorporate measurement
    void update(const MeasVec& z) {
        // Innovation
        MeasVec y = z - C * x;
        // Innovation covariance
        MeasCov S = C * P * C.T() + R;
        // Kalman gain
        GainMat K = P * C.T() * S.inverse();
        // State update
        x = x + K * y;
        // Covariance update (Joseph form for numerical stability)
        StateMat I_KC = StateMat::identity() - K * C;
        P = I_KC * P * I_KC.T() + K * R * K.T();
    }

    /// Combined predict + update
    void step(const InputVec& u, const MeasVec& z) {
        predict(u);
        update(z);
    }
};
\end{verbatim}

\textbf{Usage --- IMU tilt estimation (2-state Kalman from the case study):}
\begin{verbatim}
    KalmanFilter<2, 1, 1> kf;
    float dt = 0.001f;

    // State: [angle, gyro_bias]
    kf.A(0,0) = 1.0f;  kf.A(0,1) = -dt;
    kf.A(1,0) = 0.0f;  kf.A(1,1) = 1.0f;

    kf.B(0,0) = dt;
    kf.B(1,0) = 0.0f;

    kf.C(0,0) = 1.0f;  kf.C(0,1) = 0.0f;

    // Tune these
    kf.Q(0,0) = 0.001f;  kf.Q(1,1) = 0.003f;
    kf.R(0,0) = 0.03f;

    kf.P = Mat<2,2>::identity() * 1.0f;

    // In ISR:
    Mat<1,1> u;  u(0,0) = gyro_rate;
    Mat<1,1> z;  z(0,0) = accel_angle;
    kf.step(u, z);
    float estimated_angle = kf.x(0, 0);
    float estimated_bias  = kf.x(1, 0);
\end{verbatim}

\subsection{LQR Gain (Offline Computation Helper)}

LQR gain computation requires solving the Riccati equation, which is typically done \textit{offline} in MATLAB/Python. This module applies the precomputed gain $K$ at runtime.

\begin{verbatim}
// lqr.hpp
#pragma once
#include "mat.hpp"

template<int N_STATE, int N_INPUT>
class LQRController {
public:
    using StateVec = Mat<N_STATE, 1>;
    using InputVec = Mat<N_INPUT, 1>;
    using GainMat  = Mat<N_INPUT, N_STATE>;

    GainMat K;  // Set this from MATLAB: K = lqr(A, B, Q, R)

    /// Compute optimal input: u = -K * (x - x_ref)
    InputVec compute(const StateVec& x, const StateVec& x_ref) const {
        StateVec error = x - x_ref;
        return (K * error) * (-1.0f);
    }

    /// Compute optimal input for regulation (x_ref = 0)
    InputVec compute(const StateVec& x) const {
        return (K * x) * (-1.0f);
    }
};
\end{verbatim}

\textbf{Usage --- balance bot:}
\begin{verbatim}
    // Gain computed offline: K = lqr(A, B, diag([100,1,10,1]), 1)
    LQRController<4, 1> lqr;
    lqr.K(0,0) = -31.62f;  // theta gain
    lqr.K(0,1) = -5.27f;   // theta_dot gain
    lqr.K(0,2) = -3.16f;   // x gain
    lqr.K(0,3) = -4.12f;   // x_dot gain

    // In ISR (state from Kalman filter):
    Mat<4,1> state;
    state(0,0) = kf.x(0,0);  // theta
    state(1,0) = kf.x(1,0);  // theta_dot
    state(2,0) = encoder_pos;
    state(3,0) = encoder_vel;

    auto u = lqr.compute(state);
    set_motor_torque(u(0,0));
\end{verbatim}

\subsection{Cascaded PID}

Two-loop cascaded controller (outer position + inner velocity) as described in the gimbal case study.

\begin{verbatim}
// cascaded_pid.hpp
#pragma once
#include "pid.hpp"

class CascadedPID {
public:
    PID outer;   // Position loop (slower)
    PID inner;   // Velocity loop (faster)

    struct Config {
        PIDConfig outer_cfg;
        PIDConfig inner_cfg;
        float ff_gain = 0.0f;  // Velocity feedforward gain
    };

    void configure(const Config& cfg) {
        outer.configure(cfg.outer_cfg);
        inner.configure(cfg.inner_cfg);
        ff_gain_ = cfg.ff_gain;
    }

    /// Full cascaded update
    /// pos_setpoint: desired position
    /// pos_measured: actual position
    /// vel_measured: actual velocity
    /// vel_feedforward: reference velocity for feedforward (optional)
    float update(float pos_setpoint, float pos_measured,
                 float vel_measured, float vel_feedforward = 0.0f) {
        // Outer loop: position -> velocity command
        float vel_cmd = outer.update(pos_setpoint, pos_measured)
                      + ff_gain_ * vel_feedforward;

        // Inner loop: velocity -> torque/current command
        float output = inner.update(vel_cmd, vel_measured);
        return output;
    }

    void reset() { outer.reset(); inner.reset(); }

private:
    float ff_gain_ = 0.0f;
};
\end{verbatim}

\textbf{Usage --- gimbal yaw:}
\begin{verbatim}
    CascadedPID::Config cfg;
    // Outer loop (200 Hz)
    cfg.outer_cfg = {.kp=8.0f, .ki=0.0f, .kd=0.0f, .dt=0.005f,
                     .out_min=-300.0f, .out_max=300.0f};
    // Inner loop (1 kHz)
    cfg.inner_cfg = {.kp=20.0f, .ki=1.0f, .kd=0.0f, .dt=0.001f,
                     .out_min=-30000.0f, .out_max=30000.0f,
                     .integral_max=10000.0f};
    cfg.ff_gain = 1.0f;

    CascadedPID gimbal;
    gimbal.configure(cfg);

    // Inner loop ISR (1 kHz) - only update inner
    float torque = gimbal.inner.update(vel_cmd_cached, gyro_yaw);

    // Outer loop ISR (200 Hz) - update both
    float torque = gimbal.update(target_angle, encoder_yaw,
                                 gyro_yaw, ref_velocity);
\end{verbatim}

\subsection{Advanced PID (Integral Separation + PI-D + 2-DOF)}

Extends the base PID with all improvements from Section 4.5.

\begin{verbatim}
// pid_advanced.hpp
#pragma once
#include <cmath>

struct AdvancedPIDConfig {
    float kp = 0.0f, ki = 0.0f, kd = 0.0f;
    float dt = 0.001f;
    float out_min = -1e6f, out_max = 1e6f;
    float integral_max = 1e6f;
    float d_filter_N = 10.0f;     // Derivative LPF coefficient
    float setpoint_weight_b = 1.0f; // P setpoint weight (0-1)
    float setpoint_weight_c = 0.0f; // D setpoint weight (0=PI-D, 1=PID)
    float integral_sep_eps = 0.0f;  // 0 = disabled, >0 = threshold
};

class AdvancedPID {
public:
    AdvancedPID() = default;
    explicit AdvancedPID(const AdvancedPIDConfig& c) : cfg_(c) {}
    void configure(const AdvancedPIDConfig& c) { cfg_ = c; }

    float update(float setpoint, float measurement) {
        float error = setpoint - measurement;

        // --- P with setpoint weighting (2-DOF) ---
        float p_term = cfg_.kp * (cfg_.setpoint_weight_b * setpoint
                                  - measurement);

        // --- I with integral separation ---
        if (cfg_.integral_sep_eps <= 0.0f
            || fabsf(error) < cfg_.integral_sep_eps) {
            integral_ += cfg_.ki * error * cfg_.dt;
        }
        // Anti-windup: back-calculation
        integral_ += anti_windup_;
        integral_ = clamp(integral_, -cfg_.integral_max,
                          cfg_.integral_max);

        // --- D with setpoint weighting + filtered derivative ---
        float d_input = cfg_.setpoint_weight_c * setpoint - measurement;
        float d_raw = -(d_input - d_prev_) / cfg_.dt;
        float alpha = cfg_.dt * cfg_.d_filter_N
                    / (1.0f + cfg_.dt * cfg_.d_filter_N);
        d_filtered_ = alpha * d_raw + (1.0f - alpha) * d_filtered_;
        float d_term = cfg_.kd * d_filtered_;
        d_prev_ = d_input;

        // --- Output ---
        float raw = p_term + integral_ + d_term;
        float out = clamp(raw, cfg_.out_min, cfg_.out_max);

        // Back-calculation anti-windup
        anti_windup_ = (cfg_.ki != 0.0f)
            ? (out - raw) / cfg_.ki * cfg_.dt * 0.5f
            : 0.0f;
        return out;
    }

    void reset() {
        integral_ = 0; d_filtered_ = 0; d_prev_ = 0; anti_windup_ = 0;
    }

private:
    AdvancedPIDConfig cfg_;
    float integral_ = 0, d_filtered_ = 0, d_prev_ = 0, anti_windup_ = 0;
    static float clamp(float v, float lo, float hi) {
        return v < lo ? lo : (v > hi ? hi : v);
    }
};
\end{verbatim}

\textbf{Usage examples:}
\begin{verbatim}
    // Standard PID (same as pid.hpp)
    AdvancedPIDConfig cfg;
    cfg.kp=6; cfg.ki=0.3; cfg.kd=3; cfg.dt=0.001f;

    // PI-D (derivative on measurement, no kick)
    cfg.setpoint_weight_c = 0.0f;  // D on measurement only

    // I-PD (P and D both on measurement)
    cfg.setpoint_weight_b = 0.0f;  // P on measurement only
    cfg.setpoint_weight_c = 0.0f;  // D on measurement only

    // 2-DOF: soften setpoint response, keep disturbance rejection
    cfg.setpoint_weight_b = 0.7f;  // reduce P overshoot on setpoint
    cfg.setpoint_weight_c = 0.0f;  // no derivative kick

    // Integral separation: only integrate when close to target
    cfg.integral_sep_eps = 5.0f;   // degrees threshold
\end{verbatim}

\subsection{EKF for 6-DOF IMU (Pose Estimation)}

A ready-to-use EKF that estimates roll, pitch, yaw and gyro biases from a 6-axis IMU.

\begin{verbatim}
// ekf_imu.hpp — 6-DOF IMU Pose Estimation via EKF
// State: [roll, pitch, yaw, bias_gx, bias_gy, bias_gz]
// Measurement: accelerometer -> roll, pitch
// Part of: A Practical Guide to Control Theory
// License: MIT
#pragma once
#include "mat.hpp"
#include <cmath>

class EKF_IMU {
public:
    static constexpr int NX = 6;  // states
    static constexpr int NZ = 2;  // measurements (roll, pitch from accel)

    using State  = Mat<NX, 1>;
    using Cov    = Mat<NX, NX>;
    using MeasV  = Mat<NZ, 1>;
    using MeasC  = Mat<NZ, NZ>;
    using KGain  = Mat<NX, NZ>;
    using Hmtx   = Mat<NZ, NX>;

    State x;   // [roll, pitch, yaw, bx, by, bz]
    Cov   P;

    Cov   Q;   // Process noise covariance
    MeasC R;   // Measurement noise covariance

    float dt = 0.001f;  // sampling period

    void init() {
        x = State::zeros();
        P = Cov::identity() * 0.1f;
        // Defaults (tune for your IMU)
        Q = Cov::zeros();
        Q(0,0)=0.001f; Q(1,1)=0.001f; Q(2,2)=0.001f;
        Q(3,3)=0.0001f; Q(4,4)=0.0001f; Q(5,5)=0.0001f;
        R(0,0) = 0.15f; R(1,1) = 0.15f;
    }

    /// Main update. Call at IMU sample rate.
    /// gx,gy,gz: gyro readings (rad/s)
    /// ax,ay,az: accel readings (m/s^2)
    void update(float gx, float gy, float gz,
                float ax, float ay, float az) {
        // === PREDICT ===
        float phi = x(0,0), theta = x(1,0);
        float wx = gx - x(3,0), wy = gy - x(4,0), wz = gz - x(5,0);

        float sp = sinf(phi), cp = cosf(phi);
        float tt = tanf(theta), ct = cosf(theta);

        // Euler kinematics
        float phi_dot   = wx + sp*tt*wy + cp*tt*wz;
        float theta_dot = cp*wy - sp*wz;
        float psi_dot   = (sp/ct)*wy + (cp/ct)*wz;

        // Predict state
        State xp = x;
        xp(0,0) += phi_dot * dt;
        xp(1,0) += theta_dot * dt;
        xp(2,0) += psi_dot * dt;

        // Jacobian F (6x6)
        Cov F = Cov::identity();
        F(0,0) += (cp*tt*wy - sp*tt*wz)*dt;
        F(0,1)  = (sp/(ct*ct)*wy + cp/(ct*ct)*wz)*dt;
        F(0,3) = -dt;  F(0,4) = -sp*tt*dt;  F(0,5) = -cp*tt*dt;
        F(1,0)  = (-sp*wy - cp*wz)*dt;
        F(1,4) = -cp*dt;  F(1,5) = sp*dt;
        F(2,0)  = (cp/ct*wy - sp/ct*wz)*dt;
        float st = sinf(theta);
        F(2,1)  = (sp*st/(ct*ct)*wy + cp*st/(ct*ct)*wz)*dt;
        F(2,4) = -sp/ct*dt;  F(2,5) = -cp/ct*dt;

        Cov Pp = F * P * F.T() + Q;

        // === UPDATE (accel -> roll, pitch) ===
        float z_roll  = atan2f(ay, az);
        float z_pitch = atan2f(-ax, sqrtf(ay*ay + az*az));

        MeasV z;  z(0,0) = z_roll;  z(1,0) = z_pitch;
        MeasV h;  h(0,0) = xp(0,0); h(1,0) = xp(1,0);

        MeasV innov;
        innov(0,0) = wrap_pi(z(0,0) - h(0,0));
        innov(1,0) = wrap_pi(z(1,0) - h(1,0));

        // H = [1 0 0 0 0 0; 0 1 0 0 0 0]
        Hmtx H = Hmtx::zeros();
        H(0,0) = 1.0f; H(1,1) = 1.0f;

        MeasC S = H * Pp * H.T() + R;
        KGain Kk = Pp * H.T() * S.inverse();

        x = xp + Kk * innov;
        P = (Cov::identity() - Kk * H) * Pp;
    }

    float roll()  const { return x(0,0); }
    float pitch() const { return x(1,0); }
    float yaw()   const { return x(2,0); }
    float bias_x() const { return x(3,0); }
    float bias_y() const { return x(4,0); }
    float bias_z() const { return x(5,0); }

private:
    static float wrap_pi(float a) {
        while (a >  3.14159265f) a -= 6.28318530f;
        while (a < -3.14159265f) a += 6.28318530f;
        return a;
    }
};
\end{verbatim}

\textbf{Usage:}
\begin{verbatim}
    EKF_IMU ekf;
    ekf.dt = 0.002f;  // 500 Hz IMU
    ekf.init();

    // Tune for your IMU (MPU6050 example):
    ekf.Q(0,0) = 0.001f;  ekf.Q(1,1) = 0.001f;  ekf.Q(2,2) = 0.001f;
    ekf.Q(3,3) = 1e-5f;   ekf.Q(4,4) = 1e-5f;   ekf.Q(5,5) = 1e-5f;
    ekf.R(0,0) = 0.1f;    ekf.R(1,1) = 0.1f;

    // In IMU interrupt (500 Hz):
    ekf.update(gyro_x, gyro_y, gyro_z, accel_x, accel_y, accel_z);

    float roll_deg  = ekf.roll()  * 180.0f / M_PI;
    float pitch_deg = ekf.pitch() * 180.0f / M_PI;
    float yaw_deg   = ekf.yaw()   * 180.0f / M_PI;
\end{verbatim}

\subsection{TinyMPC Solver}

A minimal embedded MPC solver using the Riccati-based approach from TinyMPC (see Section 6 theory). Solves box-constrained LQR problems in real time on microcontrollers.

\begin{verbatim}
// tiny_mpc.hpp — Riccati-based MPC for embedded systems
// Based on: TinyMPC (Nguyen et al., 2024)
// Precomputes Riccati recursion offline, solves ADMM online.
#pragma once
#include "mat.hpp"

template<int NX, int NU, int N_HORIZON>
class TinyMPC {
public:
    using StateVec = Mat<NX, 1>;
    using InputVec = Mat<NU, 1>;
    using StateMat = Mat<NX, NX>;
    using InputMat = Mat<NX, NU>;
    using GainMat  = Mat<NU, NX>;

    // System model (discrete)
    StateMat Ad;
    InputMat Bd;

    // Cost matrices
    StateMat Q;          // State cost
    Mat<NU,NU> R;        // Input cost
    StateMat Q_terminal; // Terminal cost (set to DARE solution)

    // Constraints
    InputVec u_min, u_max;
    StateVec x_min, x_max;

    // ADMM parameters
    float rho = 1.0f;        // Penalty parameter
    int   max_iter = 10;      // ADMM iterations (5-20 typical)

    // Precomputed Riccati gains (call precompute() once offline)
    GainMat  K_gains[N_HORIZON];    // Feedback gains
    StateMat P_mats[N_HORIZON + 1]; // Cost-to-go matrices

    /// Precompute Riccati recursion (call once, or when model changes)
    void precompute() {
        P_mats[N_HORIZON] = Q_terminal;
        for (int k = N_HORIZON - 1; k >= 0; --k) {
            // P = Q + A'*P_next*A - A'*P_next*B*(R+B'*P_next*B)^-1
            //     * B'*P_next*A
            auto P_next = P_mats[k + 1];
            auto BtP = Bd.T() * P_next;
            auto S = R + BtP * Bd;  // NU x NU
            auto S_inv = S.inverse();
            K_gains[k] = S_inv * BtP * Ad;  // NU x NX
            auto AtP = Ad.T() * P_next;
            P_mats[k] = Q + AtP * Ad - AtP * Bd * K_gains[k];
        }
    }

    /// Solve MPC for current state. Returns optimal first input.
    InputVec solve(const StateVec& x0,
                   const StateVec x_ref[N_HORIZON]) {
        // Initialize with unconstrained LQR rollout
        StateVec x_traj[N_HORIZON + 1];
        InputVec u_traj[N_HORIZON];
        x_traj[0] = x0;

        for (int k = 0; k < N_HORIZON; ++k) {
            StateVec dx = x_traj[k] - x_ref[k];
            u_traj[k] = (K_gains[k] * dx) * (-1.0f);
            // Clamp (projection step)
            for (int i = 0; i < NU; ++i) {
                float v = u_traj[k](i, 0);
                v = v < u_min(i,0) ? u_min(i,0) :
                    (v > u_max(i,0) ? u_max(i,0) : v);
                u_traj[k](i, 0) = v;
            }
            x_traj[k + 1] = Ad * x_traj[k] + Bd * u_traj[k];
        }

        // ADMM refinement for constraint satisfaction
        InputVec z[N_HORIZON];   // Slack variables
        InputVec lambda[N_HORIZON]; // Dual variables
        for (int k = 0; k < N_HORIZON; ++k) {
            z[k] = u_traj[k];
            lambda[k] = InputVec::zeros();
        }

        for (int iter = 0; iter < max_iter; ++iter) {
            // Forward pass: solve unconstrained with penalty
            x_traj[0] = x0;
            for (int k = 0; k < N_HORIZON; ++k) {
                StateVec dx = x_traj[k] - x_ref[k];
                // Modified Riccati includes ADMM penalty
                InputVec u_unc = (K_gains[k] * dx) * (-1.0f);
                // ADMM correction
                u_traj[k] = u_unc + (z[k] - lambda[k]) * (rho * 0.5f);
                x_traj[k+1] = Ad * x_traj[k] + Bd * u_traj[k];
            }
            // z-update: projection onto constraints
            for (int k = 0; k < N_HORIZON; ++k) {
                for (int i = 0; i < NU; ++i) {
                    float v = u_traj[k](i,0) + lambda[k](i,0);
                    v = v < u_min(i,0) ? u_min(i,0) :
                        (v > u_max(i,0) ? u_max(i,0) : v);
                    z[k](i, 0) = v;
                }
            }
            // Dual update
            for (int k = 0; k < N_HORIZON; ++k) {
                lambda[k] = lambda[k] + u_traj[k] - z[k];
            }
        }
        return u_traj[0];
    }
};
\end{verbatim}

\textbf{Usage --- chassis trajectory tracking:}
\begin{verbatim}
    // 2D position control: state=[x,y,vx,vy], input=[ax,ay]
    TinyMPC<4, 2, 10> mpc;  // 4 states, 2 inputs, horizon=10
    float dt = 0.01f;

    // Discrete double-integrator model
    mpc.Ad = Mat<4,4>::identity();
    mpc.Ad(0,2) = dt; mpc.Ad(1,3) = dt;
    mpc.Bd(2,0) = dt; mpc.Bd(3,1) = dt;

    // Costs
    mpc.Q = Mat<4,4>::identity() * 10.0f;
    mpc.R = Mat<2,2>::identity() * 1.0f;
    mpc.Q_terminal = Mat<4,4>::identity() * 100.0f;

    // Constraints: max acceleration 5 m/s^2
    mpc.u_min(0,0) = -5; mpc.u_min(1,0) = -5;
    mpc.u_max(0,0) =  5; mpc.u_max(1,0) =  5;

    mpc.precompute();  // Once at startup

    // In control loop:
    Mat<4,1> x_ref[10];  // Fill from path planner
    auto u = mpc.solve(current_state, x_ref);
    set_accel(u(0,0), u(1,0));
\end{verbatim}

\subsection{Cubic B\'{e}zier Trajectory}

Generates smooth 2D paths through four control points. Useful for swerve drive paths and end-effector trajectories.

\begin{verbatim}
// bezier.hpp -- Cubic Bezier curve for 2D trajectory generation
#pragma once

struct Vec2 { float x, y; };

class CubicBezier {
public:
    Vec2 p0, p1, p2, p3;  // Control points

    /// Evaluate position at parameter t in [0, 1]
    Vec2 position(float t) const {
        float u = 1.0f - t;
        float u2 = u * u, t2 = t * t;
        float u3 = u2 * u, t3 = t2 * t;
        return {
            u3*p0.x + 3*u2*t*p1.x + 3*u*t2*p2.x + t3*p3.x,
            u3*p0.y + 3*u2*t*p1.y + 3*u*t2*p2.y + t3*p3.y
        };
    }

    /// Evaluate first derivative (tangent) at parameter t
    Vec2 tangent(float t) const {
        float u = 1.0f - t;
        return {
            3*(u*u*(p1.x-p0.x) + 2*u*t*(p2.x-p1.x) + t*t*(p3.x-p2.x)),
            3*(u*u*(p1.y-p0.y) + 2*u*t*(p2.y-p1.y) + t*t*(p3.y-p2.y))
        };
    }

    /// Evaluate second derivative (curvature direction)
    Vec2 acceleration(float t) const {
        float u = 1.0f - t;
        return {
            6*(u*(p2.x - 2*p1.x + p0.x) + t*(p3.x - 2*p2.x + p1.x)),
            6*(u*(p2.y - 2*p1.y + p0.y) + t*(p3.y - 2*p2.y + p1.y))
        };
    }

    /// Approximate arc length by summing N linear segments
    float arc_length(int N = 100) const {
        float len = 0;
        Vec2 prev = p0;
        for (int i = 1; i <= N; ++i) {
            Vec2 cur = position(static_cast<float>(i) / N);
            float dx = cur.x - prev.x, dy = cur.y - prev.y;
            len += sqrtf(dx*dx + dy*dy);
            prev = cur;
        }
        return len;
    }
};
\end{verbatim}

\textbf{Usage --- swerve drive path:}
\begin{verbatim}
    CubicBezier path;
    path.p0 = {0, 0};        // Start
    path.p1 = {1.0f, 0};     // Exit direction: forward
    path.p2 = {2.0f, 1.5f};  // Entry direction at end
    path.p3 = {3.0f, 1.5f};  // End point

    // Track with 50 samples
    for (int i = 0; i <= 50; ++i) {
        float t = static_cast<float>(i) / 50;
        Vec2 pos = path.position(t);
        Vec2 vel = path.tangent(t);  // For feedforward heading
        // Send (pos.x, pos.y) to chassis controller
    }
\end{verbatim}

\subsection{Module Dependency Map}

\begin{center}
\begin{tabular}{lll}
\toprule
\textbf{Module} & \textbf{File} & \textbf{Depends on} \\
\midrule
Low-pass filter & \texttt{lpf.hpp} & (none) \\
PID controller & \texttt{pid.hpp} & (none) \\
Advanced PID & \texttt{pid\_advanced.hpp} & (none) \\
Setpoint ramp & \texttt{ramp.hpp} & \texttt{lpf.hpp} \\
Complementary filter & \texttt{complementary\_filter.hpp} & (none) \\
Matrix library & \texttt{mat.hpp} & (none) \\
Kalman filter & \texttt{kalman.hpp} & \texttt{mat.hpp} \\
EKF IMU (6-DOF) & \texttt{ekf\_imu.hpp} & \texttt{mat.hpp} \\
LQR controller & \texttt{lqr.hpp} & \texttt{mat.hpp} \\
TinyMPC solver & \texttt{tiny\_mpc.hpp} & \texttt{mat.hpp} \\
Cascaded PID & \texttt{cascaded\_pid.hpp} & \texttt{pid.hpp} \\
Cubic B\'{e}zier & \texttt{bezier.hpp} & (none) \\
\bottomrule
\end{tabular}
\end{center}

All modules are MIT licensed. Copy into your project's \texttt{include/control/} directory, \texttt{\#include}, and go.

\newpage

\vspace{1em}
\hrule
\vspace{1em}

\section*{Closing Thoughts}

If you have made it this far, congratulations---you now have a solid theoretical foundation for designing control systems on competition robots. Here is a cheat sheet for what to reach for in different situations:

\begin{center}
\begin{tabular}{ll}
\toprule
\textbf{Situation} & \textbf{Tool} \\
\midrule
Noisy sensor data & Low-pass filter (Section 2) \\
What does my motor actually do? & Motor modeling (Section 3) \\
Single motor speed/position control & PID (Section 4) \\
Gimbal tracking & PID + feedforward (Section 4) \\
Smooth point-to-point moves & Trajectory generation (Section 6) \\
Getting PID from paper to MCU & Tustin discretization (Section 5) \\
Balance bot / MIMO system & LQR (Section 6) \\
Trajectory tracking with constraints & MPC (Section 6) \\
State estimation from noisy sensors & Kalman filter (Section 6) \\
IMU + encoder fusion & Extended Kalman filter (Section 6) \\
Robot can flip/tumble past $\pm 90$\textdegree{} & Quaternion EKF (Section 6) \\
\bottomrule
\end{tabular}
\end{center}

Remember: theory is a tool, not an end goal. The best controller is the one that \textbf{works on the field}, not the one that looks most elegant on paper. Start simple (PID), understand your system, add complexity only when needed, and always, always test on real hardware.

Good luck at your next competition. Go build something awesome.

